\section{讨论与误差分析}

\subsection{主要误差来源}

\textbf{遮挡。} 一个空间点可能只在左图或右图可见。匹配器在遮挡边界无法找到真正
的对应点，容易产生空洞、孤立视差或边界附近的跳变；有效掩码中的空白区域应当保留，
而不是通过插值伪造测量结果。

\textbf{纹理不足。} 纯色或低对比度区域内，多个候选视差的匹配代价相近，
\texttt{uniquenessRatio} 会主动拒绝不够明确的结果。桌面和背景中的大块弱纹理区域因此更
容易出现无效值或不稳定视差。

\textbf{重复纹理。} 周期性纹理会产生多个近似相同的局部匹配代价。SGBM 的路径平滑
可以利用邻域结构减少部分歧义，但不能保证在重复图案上找到正确的全局对应关系。

\textbf{图像边界与搜索范围。} 匹配窗口和视差搜索范围会消耗左右边缘的可用像素。
当真实视差超出 \texttt{minDisparity} 到 \texttt{minDisparity+numDisparities} 的搜索区间时，
即使图像中存在纹理，也无法得到有效结果。

\subsection{参数和下采样的影响}

本次运行只记录了一组固定参数，没有进行额外对比实验。一般而言，增大 \texttt{blockSize}
可以提高弱纹理区域的局部稳定性，但会使物体边界变粗；减小窗口有利于保留细节，却
更容易受到噪声影响。\texttt{P1} 和 \texttt{P2} 控制视差平滑程度，过大可能抹平真实边界，过小则
可能产生散斑。\texttt{uniquenessRatio} 提高后会留下更少但更可靠的匹配；speckle 参数则
用于删除小面积的异常连通区域。

同步执行一次 \texttt{pyrDown} 可以将匹配像素数和计算负担显著降低，适合课程实验的快速
复现；代价是细小叶片边缘和小纹理被合并，视差细节减少。若以细节为优先，应在保持
左右缩放一致的前提下关闭降采样，并重新记录一整套数组、图片和参数，而不能把不同
分辨率的结果混在同一张表中。

\subsection{深度尺度限制}

当前结果只测量像素视差。根据 $Z=fB/d$，即使视差估计完全正确，未知的 $fB$ 仍会使
深度存在任意比例因子；因此 \codepath{depth_proxy} 的数值不能解释为米或厘米。官方样例目录
中的其他标定文件没有被证明与 aloe 图像、分辨率和相机设置配套，不能直接套用。

\subsection{有效像素比例不等于匹配准确率}

本次运行得到的 $73.4713\%$ 有效像素比例，严格含义是有多少处理分辨率中的像素通过了
“有限且 $d>15$”这一筛选。它可以反映结果的覆盖范围，却不能回答这些视差是否与真实
对应点一致。一个像素即使获得了有限视差，也可能来自重复纹理、遮挡边界或错误的局部
代价最小值；反过来，无效像素也可能只是因为匹配器无法给出足够有把握的答案，而不代表
该处场景不存在深度。因此，本实验把覆盖率、视差分布和匹配正确性分开讨论，没有把
有效比例写成准确率或识别率。

若要得到可比较的准确性指标，需要额外提供与当前分辨率和坐标定义一致的视差参考，或
使用独立的左右一致性误差、重投影误差和人工标注区域进行评价。即使保留官方
\texttt{aloeGT.png}，也应先核对它的编码尺度、有效区域和与当前图像的几何对应关系，
再计算 MAE、RMSE 或 bad-pixel ratio；在这些条件未被验证前，不能把参考图直接当作本实验
的真值。当前报告的结论因此限定为“产生了可解释的相对视差结构”，而不是“达到了某个
视差准确率”。

\subsection{从数值分布到空间结构的解释}

视差范围和分位数能够说明数值是否落在匹配器的工作区间，但不能单独说明误差发生在
哪里。空间上，遮挡边界、弱纹理区域以及图像两侧的搜索边界可能同时表现为无效像素；
相同的无效比例也可能对应完全不同的空间分布。因而，分析结果时应把有效掩码作为一张
空间支持图，并与左图纹理、视差灰度和相对深度图并排观察。这样可以区分“整块区域没有
足够证据”和“物体边界附近少量像素被拒绝”这两类情况，避免只用一个总体百分比概括
所有失效模式。

同理，视差较大的区域通常可解释为相对更近，但单个高视差斑块也可能是重复纹理造成的
错误匹配。SGBM 的平滑项会利用邻域连续性补足局部证据，因此连续的结果既可能反映真实
表面，也可能反映算法先验。只有当空间结构、左右一致性和可用参考共同支持时，才适合
把局部视差变化解释为物体前后关系。

\subsection{单像素扰动对远近区域的不同影响}

由式~\eqref{eq:depth-sensitivity} 可知，深度反演是非线性的。对相同的绝对视差误差
$\Delta d$，小视差区域的相对误差幅度更大；这通常对应较远的场景区域。因而，远处弱
纹理区域即使只出现约一个像素量级的匹配偏差，也可能在相对深度排序中造成更明显的
扰动。近处区域具有较大的视差，前后关系通常更容易分开，但遮挡和边界混合会增加局部
匹配困难。这说明“远处的数值更稳定”或“近处一定更准确”都不是由公式自动保证的，
评价时应同时考虑视差量级、纹理条件和边界位置。

该非线性也是本实验保留浮点视差而不只保存 8-bit 灰度图的原因。显示图经过第 1 至第
99 百分位拉伸后，主要用于观察结构；它会压缩极端值并丢失无效值的数值语义。任何关于
视差扰动、深度变化或分位数的分析，都应回到 \texttt{disparity\_raw.npy} 和
\texttt{depth\_proxy.npy}，而不能从显示灰度反推原始像素位移。

\subsection{低成本改进方案}

在不改变课程任务范围的前提下，可以采用以下低成本改进：

\begin{enumerate}
  \item 对同一双目系统采集带已知方格尺寸的左右棋盘格序列，估计内参、畸变、外参和 $Q$；
  \item 使用标定参数先完成去畸变和立体校正，再运行同一套 SGBM 与数组保存流程；
  \item 在保留原始 \texttt{npy} 的前提下，对 \texttt{numDisparities}、\texttt{blockSize} 和 speckle 参数做少量
        预先规定的对比，并把每组参数与结果分开保存；
  \item 对视差结果增加左右一致性和置信度统计，用于区分“无匹配”和“低置信匹配”。
\end{enumerate}

这些改进的优先级并不相同：若目标是物理测量，应先完成同一相机系统的标定和校正；
若目标只是比较匹配器的空间结构，则应先增加预先规定的参数对照和一致性统计。两种
目标都应保持输入、分辨率、搜索范围和输出文件一一对应，避免把不同实验条件下的
视差数组混入同一组结论。以上是改进建议，本次提交没有声称已经完成这些额外标定或
对比实验。
